\subsubsection{金属平均族的闭式计数、有限尺度振荡与整数尺度相变}\label{subsubsec:ostrowski-metallic-binet-oscillation-integer-phase-transition}
在命题 \ref{prop:metallic-two-state-sft}、命题 \ref{prop:metallic-compression-extremum} 与命题 \ref{prop:metallic-linear-scalarization-threshold} 已经把常数型金属平均族的两态语法、连续压缩极值与连续标量化前沿固定下来之后，有限尺度计数与整数尺度选择还允许进一步闭式化。
下述结果把合法语言计数、有限分辨率压缩缺陷与离散样本阈值统一到同一组完全显式公式中。

\begin{theorem}[金属平均族稳定语言计数的闭式 Binet 公式]\label{thm:metallic-binet-closed-form}
沿命题 \ref{prop:metallic-two-state-sft} 的记号，固定整数 \(A\ge 1\)，并记
\[
\alpha_A=[0;A,A,A,\dots],
\qquad
\lambda_A=\frac{A+\sqrt{A^2+4}}{2},
\qquad
\mu_A:=\frac{A-\sqrt{A^2+4}}{2}=-\lambda_A^{-1}.
\]
令 \((q_n)_{n\ge 0}\) 为 \(\alpha_A\) 的收敛分母序列，即
\[
q_{-1}=0,\qquad q_0=1,\qquad q_{n+1}=Aq_n+q_{n-1}\quad(n\ge 0).
\]
则对一切 \(n\ge 0\) 都有严格闭式
\[
q_n=
\frac{\lambda_A^{n+2}+(-1)^n\lambda_A^{-n}}{\lambda_A^2+1}.
\]
特别地，由命题 \ref{prop:Xm-alpha-cardinality}，对一切 \(m\ge 1\),
\[
\bigl|X_m^{(\alpha_A)}\bigr|
=q_{m+1}
=
\frac{\lambda_A^{m+3}+(-1)^{m+1}\lambda_A^{-m-1}}{\lambda_A^2+1}.
\]
\leanverified{paper\_metallic\_binet\_closed\_form}
\leanverified{phi\_irrational}
\end{theorem}
\begin{proof}
递推 \(q_{n+1}=Aq_n+q_{n-1}\) 的特征方程为
\[
\lambda^2-A\lambda-1=0,
\]
其两根正是 \(\lambda_A\) 与 \(\mu_A\)。
因此存在常数 \(c_1,c_2\) 使
\[
q_n=c_1\lambda_A^n+c_2\mu_A^n.
\]
由初值 \(q_{-1}=0,\ q_0=1\) 解得
\[
q_n=\frac{\lambda_A^{n+1}-\mu_A^{n+1}}{\lambda_A-\mu_A}.
\]
再用 \(\mu_A=-\lambda_A^{-1}\) 以及
\[
\lambda_A-\mu_A
=
\lambda_A+\lambda_A^{-1}
=
\frac{\lambda_A^2+1}{\lambda_A}
\]
化简，即得
\[
q_n=
\frac{\lambda_A^{n+2}+(-1)^n\lambda_A^{-n}}{\lambda_A^2+1}.
\]
最后代入命题 \ref{prop:Xm-alpha-cardinality} 的恒等式
\(
|X_m^{(\alpha_A)}|=q_{m+1}
\)
即可。证毕。
\end{proof}

\leanverified{paper\_metallic\_finite\_scale\_compression\_oscillation}
\begin{theorem}[金属平均族投影压缩率的有限尺度奇偶振荡公式]\label{thm:metallic-finite-scale-compression-oscillation}
沿用前述记号。定义有限尺度投影压缩量
\[
\Xi_m(A):=\log\frac{(A+1)^m}{|X_m^{(\alpha_A)}|}
\qquad (m\ge 1),
\]
以及常数项
\[
\gamma_A:=\log\frac{\lambda_A^2+1}{\lambda_A^3}.
\]
则对一切 \(m\ge 1\) 都有精确恒等式
\[
\Xi_m(A)
=
m\,h_Z(A)+\gamma_A-\log\!\Bigl(1+(-1)^{m+1}\lambda_A^{-2m-4}\Bigr).
\]
从而对任意 \(n\ge 0\)，偶、奇子序列满足
\[
\Xi_{2n}(A)-\bigl(2n\,h_Z(A)+\gamma_A\bigr)
=
-\log\!\bigl(1-\lambda_A^{-4n-4}\bigr)>0,
\]
\[
\Xi_{2n+1}(A)-\bigl((2n+1)\,h_Z(A)+\gamma_A\bigr)
=
-\log\!\bigl(1+\lambda_A^{-4n-6}\bigr)<0.
\]
因此有限尺度误差并非仅具 \(O(\lambda_A^{-2m})\) 量级，而是满足严格的奇偶振荡律：偶数子序列从上方收敛，奇数子序列从下方收敛，且两者都以精确指数 \(2m\) 收敛到同一仿射主项。
\end{theorem}
\begin{proof}
由定理 \ref{thm:metallic-binet-closed-form}，
\[
\bigl|X_m^{(\alpha_A)}\bigr|
=
\frac{\lambda_A^{m+3}+(-1)^{m+1}\lambda_A^{-m-1}}{\lambda_A^2+1}
=
\frac{\lambda_A^{m+3}}{\lambda_A^2+1}
\Bigl(1+(-1)^{m+1}\lambda_A^{-2m-4}\Bigr).
\]
因此
\[
\Xi_m(A)
=
m\log(A+1)
-(m+3)\log\lambda_A
+\log(\lambda_A^2+1)
-\log\!\Bigl(1+(-1)^{m+1}\lambda_A^{-2m-4}\Bigr).
\]
把前两项合并为
\[
m\bigl(\log(A+1)-\log\lambda_A\bigr)=m\,h_Z(A),
\]
余下常数项正是
\[
\log(\lambda_A^2+1)-3\log\lambda_A
=
\log\frac{\lambda_A^2+1}{\lambda_A^3}
=
\gamma_A.
\]
于是得到主公式。
再分别代入 \(m=2n\) 与 \(m=2n+1\) 即得两条奇偶公式，其符号由 \(\lambda_A>1\) 直接给出。证毕。
\end{proof}

\begin{theorem}[整数尺度的线性标量化阈值]\label{thm:metallic-integer-scalarization-threshold}
沿用命题 \ref{prop:metallic-linear-scalarization-threshold} 的标量化目标
\[
J_\beta(A):=h_Z(A)-\beta\,\kappa(A),
\qquad
\kappa(A)=\frac{A}{\log\lambda_A}.
\]
在整数约束 \(A\in\ZZ_{\ge 1}\) 下，帕累托集合缩减为 \(\{1,2\}\)。
因此存在唯一阈值
\[
\beta_{\ZZ}
:=
\frac{h_Z(2)-h_Z(1)}{\kappa(2)-\kappa(1)}
=
\frac{\log\!\left(\dfrac{3\varphi}{2(1+\sqrt2)}\right)}
{\dfrac{2}{\log(1+\sqrt2)}-\dfrac{1}{\log\varphi}}
\approx 0.0277519138993265,
\]
其中 \(\varphi=(1+\sqrt5)/2\)。
并且最优整数尺度满足严格相变律
\[
A_\beta^{\ZZ}
=
\begin{cases}
2,& 0\le \beta<\beta_{\ZZ},\\[2mm]
\{1,2\},& \beta=\beta_{\ZZ},\\[2mm]
1,& \beta>\beta_{\ZZ}.
\end{cases}
\]
\end{theorem}
\begin{proof}
由注 \ref{rem:metallic-design-space-reduction}，所有整数 \(A\ge 3\) 在“最大化 \(h_Z\)、最小化 \(\kappa\)”的双目标意义下都被 \(A=2\) 严格支配。
故在线性标量化
\(
J_\beta(A)=h_Z(A)-\beta\kappa(A)
\)
下，只需比较 \(A=1\) 与 \(A=2\)。

由 \(\lambda_1=\varphi\) 与 \(\lambda_2=1+\sqrt2\)，有
\[
h_Z(1)=\log\frac{2}{\varphi},
\qquad
h_Z(2)=\log\frac{3}{1+\sqrt2},
\]
\[
\kappa(1)=\frac{1}{\log\varphi},
\qquad
\kappa(2)=\frac{2}{\log(1+\sqrt2)}.
\]
阈值条件 \(J_\beta(1)=J_\beta(2)\) 唯一给出
\[
\beta
=
\frac{h_Z(2)-h_Z(1)}{\kappa(2)-\kappa(1)}
=
\beta_{\ZZ}.
\]
由于 \(\kappa(2)>\kappa(1)\) 且 \(h_Z(2)>h_Z(1)\)，当 \(\beta<\beta_{\ZZ}\) 时压缩项占优，故 \(A=2\) 唯一胜出；当 \(\beta>\beta_{\ZZ}\) 时证书项占优，故 \(A=1\) 唯一胜出；等号时二者并列。证毕。
\end{proof}
\leanverified{paper_metallic_integer_scalarization_threshold}

\begin{corollary}[样本惩罚律下的整数相变与 Lambert--\texorpdfstring{$W$}{W} 闭式]\label{cor:metallic-sample-penalty-integer-threshold-lambertw}
\leanverified{paper\_metallic\_sample\_penalty\_integer\_threshold\_lambertw}
沿用定理 \ref{thm:metallic-integer-scalarization-threshold} 的阈值 \(\beta_{\ZZ}\)。
对整数样本量 \(N\ge 3\)，令
\[
\beta(N):=\frac{\log N}{N},
\]
并按同一标量化目标在整数尺度上取最优点 \(A_N^{\ZZ}\)。
则存在唯一实阈值
\[
N_{\ZZ}
=
-\frac{W_{-1}(-\beta_{\ZZ})}{\beta_{\ZZ}}
\approx 188.850182986328,
\]
其中 \(W_{-1}\) 为 Lambert--\(W\) 函数的 \((-1)\) 支。
于是对所有整数 \(N\ge 3\)，最优整数尺度满足
\[
A_N^{\ZZ}
=
\begin{cases}
1,& 3\le N\le 188,\\[2mm]
2,& N\ge 189.
\end{cases}
\]
\end{corollary}
\begin{proof}
由定理 \ref{thm:metallic-integer-scalarization-threshold}，整数最优尺度完全由 \(\beta(N)\) 与 \(\beta_{\ZZ}\) 的比较决定。
函数
\[
\beta(N)=\frac{\log N}{N}
\]
在 \(N\ge e\) 上严格递减，而 \(3>e\)，故其在全部整数 \(N\ge 3\) 上也严格递减。
阈值点满足方程
\[
\frac{\log N}{N}=\beta_{\ZZ}.
\]
令 \(u=\log N\)，则 \(N=e^u\)，上式等价于
\[
u\,e^{-u}=\beta_{\ZZ}
\qquad\Longleftrightarrow\qquad
(-u)e^{-u}=-\beta_{\ZZ}.
\]
因此
\[
u=-W(-\beta_{\ZZ}).
\]
为得到 \(N>e\) 的大根，必须选取 Lambert--\(W\) 的 \((-1)\) 支，从而
\[
N_{\ZZ}=e^u
=
e^{-W_{-1}(-\beta_{\ZZ})}
=
-\frac{W_{-1}(-\beta_{\ZZ})}{\beta_{\ZZ}}.
\]
数值上 \(N_{\ZZ}\approx 188.85018\)。
由于 \(\beta(N)\) 在 \(N\ge 3\) 上严格递减，故
\[
\beta(N)>\beta_{\ZZ}\iff N\le 188,
\qquad
\beta(N)<\beta_{\ZZ}\iff N\ge 189.
\]
再由定理 \ref{thm:metallic-integer-scalarization-threshold} 的相变律得到结论。证毕。
\end{proof}

\endinput
